\documentclass[exercice,a4]{thedocument}%
\usepackage[french]{babel}%
\usepackage{array}%
\usepackage[np]{numprint}%
\usepackage{graphicx}%
\usepackage{multicol}%
\usepackage{amsmath}%
\startdatakeys
    \source{22GENMATMEAG1}
    \BO{2026}
    \cycle{4}
    \competencesDuSocle{CA,RA}
    \objectifsApprentissage{C4C2aa,C4D2ae,C4B1bf}
    \motsCles{volume,cylindre,conversion,pourcentage,proportionnalité}
\enddatakeys
\begin{document}%

\medskip

Dans une habitation, la consommation d'eau peut être anormalement élevée lorsqu’il y a une fuite d'eau.

On considère la situation suivante :
\begin{itemize}
\item Une salle de bain est équipée d'une vasque de forme cylindrique, comme l'illustre l'image ci-dessous.

\item Le robinet fuit à raison d'une goutte par seconde.

\item En moyenne, 20 gouttes d'eau correspondent à un millilitre (\np[ml]{1} ).
\end{itemize}

\begin{multicols}{2}
\begin{minipage}[7cm]{8cm}
\phantom{x}
\ps[scale=0.5]{vasque-2}
\end{minipage}

\columnbreak

~\\\\\\

\begin{tabular}{|l|}\hline

Caractéristiques de la vasque :\\

\hspace{2em}Diamètre intérieur : \np[cm]{40}\\

\hspace{2em}Hauteur intérieure : \np[cm]{15} \\

\hspace{2em}Masse : \np[kg]{25} \\\hline
\end{tabular}
\end{multicols}
\medskip

Rappels : \hspace{4em} \begin{tabular}{|c|}\hline

Volume du cylindre = $\pi \times \text{rayon}^2 \times \text{hauteur}$\\

\hspace{2.5em} \np[dm^3]{1} = 1 litre\\
\hline
\end{tabular}

\begin{enumerate}
\item En raison de la fuite, montrer qu'il tombe \np{86400}~gouttes dans la vasque en une journée complète.

\item Calculer, en litres, le volume d'eau qui tombe dans la vasque en une semaine en raison de la fuite.

\item  Montrer que la vasque  a un volume de \np{18.85}~litres, arrondi au centilitre près.

\item L'évacuation de la vasque est fermée et le logement inoccupé pendant une semaine. L'eau va-t-elle déborder de la vasque ? Justifier la réponse.
\item  À la fin du XIX\ieme{} siècle, la consommation domestique d'eau par habitant en France était d'environ 17 litres par jour. Elle a fortement augmenté avec la généralisation de la distribution d'eau par le robinet dans les domiciles : elle est passée à $165$~litres par jour et par habitant en 2004.

En 2018, la consommation des Français baisse légèrement pour atteindre $148$~litres d'eau par jour et par habitant.

Calculer le pourcentage de diminution de la consommation quotidienne d'eau par habitant entre 2004 et 2018. On arrondira ce pourcentage à l'unité.
\end{enumerate}

\showthesource
\end{document}
